% generated by roam.py handbook export -- не редактировать, править заметки
% 20261025161604-квазиньютоновские_методы.org (body deb031ca6fc4)

\begin{method}[{Квазиньютоновские методы}]\label[method]{note:6BA18878-DE34-46A5-ABE2-EDA09CAE56A7}
  Основаны на \hyperref[note:C0ED5501-CFC6-42BA-A53A-1B346CBA545D]{методе Ньютона}.
  Они призваны сохранить его плюсы. Но при этом минимизировать его минусы, как обращение матрицы (которая затратна по скорости вычисления) и хранение \hyperref[note:592AD816-1D80-48E6-8399-8524FBC1E8F3]{гессиана} (много памяти нужно)
\end{method}

\paragraph*{Общая идея}
Тут мы используем последовательность матриц, а не просто матрицу вторых производных.
Мы как бы стараемся приблизить вторую производную этими матрицами.
\[
x_{k+1} = x_k - B_k^{-1} \nabla f(x_k),\quad f(x) \approx f(x_k) + \left\langle \nabla f(x_k); x -x_k \right\rangle + \frac{1}{2} \left\langle B_k(x-x_k); x-x_k \right\rangle
\]

Хотим, чтобы $B_{k+1} = F(B_k, ...) \Rightarrow B_{k+1}^{-1} = G (B_k^{-1}, ...)$

Какие условия накладываются на $B_k$?
\[
B_k = B_k^T\quad B_k > 0\quad y_k = B_{k+1} d_k
\]
где $y_k = B_{k+1}d_k$~--- \hyperref[note:E976DE49-0091-45E1-B621-A236F8B87859]{Условие секущей}

Рассмотрим $B_k$ как неизвестную. Ее задают $\frac{n(n+1)}{2}$ неизвестных скаляров.
Посчитаем сколько уравнений на эти параметры мы получим из наших условий:

\begin{itemize}
  \item Условие положительности дает $n$ неравенств по критерию Сильвестра
  \item Условие секущей дает тоже $n$ уравнений
\end{itemize}

Получается неполная система. Каждый метод сам определяет, как дополнить эту систему.
Отсюда и выходит разнообразие методов.

\paragraph*{Методы}
\begin{itemize}
  \item \hyperref[note:01306845-D37F-480A-8DCE-EE9CA022DA1A]{BFGS}
  \item \hyperref[note:F7506DC4-37FB-4E88-8178-B0FC678233C0]{SR1}
\end{itemize}
